\documentclass[11pt]{article}
\usepackage[margin=1in]{geometry}
\usepackage{amsmath,amssymb}
\usepackage{array}
\usepackage{booktabs}
\usepackage{hyperref}

\title{M.I.N.D.: Reliability-Gated Motor Imagery EEG Decoding with Calibrated Classical Ensembles}
\author{JR-Academy 6857}
\date{}

\begin{document}
\maketitle

\begin{abstract}
Motor-imagery brain-computer interfaces (BCIs) are often evaluated as ordinary
classifiers, but assistive control has a stricter requirement: an uncertain
decoder should withhold action rather than issue a wrong command. This paper
presents M.I.N.D., a reliability-focused motor-imagery EEG decoding pipeline
that combines label-free session adaptation, filter-bank common spatial
patterns (FBCSP), calibrated LDA and SVM classifiers, equal-weight probability
voting, and confidence-based selective prediction. On BCI Competition IV
Dataset 2a, the LDA+SVM ensemble improves all-trial accuracy over LDA
(\(0.525\) vs. \(0.435\)), lowers mean Expected Calibration Error (ECE)
(\(0.141\) vs. \(0.176\)), and lowers mean multiclass Brier score
(\(0.591\) vs. \(0.672\)). At a matched 60\% operating coverage, the ensemble
also improves selective accuracy over LDA by \(0.089\) on average across the
nine subjects. A small MOABB external-validation subset using Cho2017 and
PhysionetMI shows that matched-coverage accuracy and Brier-score improvements
persist beyond BCI IV 2a, although ECE worsens in that tiny subset. This
writeup is retained as the classical baseline record for the final paper;
modern deep models, IIIa validation, broader MOABB validation, calibration
ablations, and controller replay belong to the v2 paper branch. The central
contribution is not a benchmark-accuracy claim; it is a reproducible
framework for evaluating whether a BCI decoder is reliable enough to act.
\end{abstract}

\section{Purpose}

Motor-imagery BCIs decode imagined movement from EEG. The long-term assistive
goal is compelling: a person with impaired movement could use neural activity
to control a prosthetic or communication device. The practical problem is that
EEG is low-amplitude, noisy, and non-stationary. A model that looks acceptable
under average classification accuracy can still be unusable if it fires
incorrect commands at the wrong time.

M.I.N.D. is built around a control-safety principle:
\[
  \text{a wrong action is worse than no action.}
\]
Therefore, the project evaluates the decoder as a selective controller. When
confidence is high, the system emits a command; when confidence is low, it
abstains. This makes reliability metrics such as coverage, risk, calibration,
and Brier score central rather than secondary.

\section{Background Knowledge}

\subsection{Motor Imagery EEG}

Motor imagery produces measurable changes in sensorimotor rhythms, especially
in the mu and beta ranges. These signals are subject-specific and vary across
recording sessions due to electrode placement, fatigue, artifacts, and changes
in cognitive strategy. This makes cross-session evaluation more realistic than
within-session cross-validation.

\subsection{FBCSP}

Common Spatial Patterns (CSP) learns spatial filters that emphasize
class-discriminative variance patterns. M.I.N.D. uses filter-bank CSP over
four bands:
\[
  8\!-\!12,\quad 12\!-\!16,\quad 16\!-\!20,\quad 20\!-\!30\ \mathrm{Hz}.
\]
Bandpower features are concatenated with CSP features. This keeps the main
pipeline interpretable and compatible with small EEG datasets.

\subsection{Calibration and Selective Prediction}

The base classifiers output class probabilities after sigmoid calibration.
For a probability vector \(p(y \mid x)\), confidence is defined as
\[
  c(x) = \max_y p(y \mid x).
\]
At threshold \(t\), the decoder fires only when \(c(x) \geq t\). We report:
\begin{itemize}
  \item coverage: the fraction of trials that fire,
  \item risk: error rate among fired trials,
  \item ECE: confidence-accuracy mismatch across confidence bins,
  \item Brier score: squared error of the full probability distribution.
\end{itemize}

\section{Architecture and Novel Application}

\subsection{System Framing}

M.I.N.D. is best understood as a three-layer assistive-control architecture
rather than a single classifier. The first layer is a signal layer that turns
raw motor-imagery EEG into subject-specific features while reducing
cross-session covariance shift. The second layer is a probabilistic decoding
layer that compares simple calibrated models and combines complementary
decision boundaries by soft voting. The third layer is a control layer that
decides whether the probability evidence is strong enough to issue an action.

This separation is important for the claim. The project does not argue that
one classifier is universally best. It argues that an assistive BCI should
expose its uncertainty to the controller, and that the controller should be
allowed to withhold action when the neural evidence is weak. In this framing,
coverage is a usability variable, selective risk is a safety variable, and
calibration is the bridge between the machine-learning model and the user's
trust in the device.

\subsection{Label-Free Session Adaptation}

The primary evaluation trains each subject-specific decoder on BCI Competition
IV 2a training session \(A0xT\) and scores once on the paired evaluation
session \(A0xE\). To reduce cross-session covariance shift without using
evaluation labels, M.I.N.D. applies Euclidean alignment independently to each
session.

Let \(X_i \in \mathbb{R}^{C \times T}\) be the \(i\)-th epoch. Its sample
covariance is
\[
  S_i = \frac{(X_i - \bar{X}_i)(X_i - \bar{X}_i)^\top}{T - 1}.
\]
The regularized covariance is
\[
  \Sigma_i = S_i + \epsilon \frac{\operatorname{tr}(S_i)}{C}I,
\]
where \(\epsilon = 10^{-6}\). The session reference covariance is
\[
  \bar{\Sigma} = \frac{1}{N}\sum_{i=1}^{N}\Sigma_i.
\]
Each epoch is aligned by
\[
  X_i^{\mathrm{EA}} = \bar{\Sigma}^{-1/2}X_i.
\]
The evaluation labels are used only after all predictions have been produced.
No \(A0xE\) labels are used to fit the alignment transform, feature extractor,
classifier, calibration model, threshold, or ensemble rule.

\subsection{Main Decoder}

After alignment, the main decoder is:
\[
  \text{EA} \rightarrow \text{FBCSP + bandpower} \rightarrow
  \text{calibrated LDA},
\]
\[
  \text{EA} \rightarrow \text{FBCSP + bandpower} \rightarrow
  \text{calibrated SVM}.
\]
The ensemble is an equal soft vote:
\[
  p_{\mathrm{ens}}(y \mid x) =
  \frac{1}{2}p_{\mathrm{LDA}}(y \mid x)
  + \frac{1}{2}p_{\mathrm{SVM}}(y \mid x).
\]
The pre-specified classical comparison is calibrated LDA. This keeps the
central research question focused: does a calibrated probability ensemble plus
abstention improve reliability over a strong, simple classical baseline?

\subsection{Decision Contract}

The decoder output is not treated as a direct prosthetic command. It is treated
as a proposed command with an attached reliability score. At deployment time,
the downstream device would receive one of two states:
\[
  \text{FIRE}(y) \quad \text{or} \quad \text{HOLD}.
\]
This command contract gives the user and device a predictable failure mode. A
low-confidence trial does not become a random movement class; it becomes an
explicit non-action. The benchmark evaluation is therefore aligned with the
intended use case: a model is better only if it improves correct actions at
matched coverage or lowers risk without hiding unusably low coverage.

\subsection{EEGNet Comparison}

EEGNet is included as a modern deep-learning comparison only. The project does
not use EEGNet inside the ensemble and does not tune the architecture for
benchmark optimization. The optional implementation follows the standard
EEGNet pattern: temporal convolution, depthwise spatial convolution,
separable temporal convolution, dropout, and a softmax classifier. It emits
probabilities into the same evaluation tables as LDA and the ensemble:
risk-coverage curves, ECE, Brier score, and reliability diagrams. This allows
the paper to compare the reliability-controlled classical pipeline against a
recognizable deep baseline without making deep learning the project's focus.

\section{Evaluation}

\subsection{Datasets}

The primary dataset is BCI Competition IV Dataset 2a, using nine subjects and
the train/evaluation session split \(A0xT \rightarrow A0xE\). External
validation is implemented through MOABB. The current external smoke subset
uses one subject each from Cho2017 and PhysionetMI; the code supports larger
MOABB runs by increasing the number of subjects.

\subsection{Metrics}

The evaluation reports standard accuracy and reliability-focused metrics:
\begin{itemize}
  \item all-trial accuracy,
  \item matched-coverage selective accuracy,
  \item risk-coverage curves,
  \item Expected Calibration Error,
  \item multiclass Brier score,
  \item reliability diagrams.
\end{itemize}
For the headline BCI IV 2a analysis, the operating coverage is fixed at
60\%. The subject-level ensemble-vs-LDA deltas are summarized with a
Wilcoxon signed-rank test.

\subsection{BCI IV 2a Findings}

\begin{center}
\begin{tabular}{lrrr}
\toprule
Metric & Ensemble & LDA & Ensemble - LDA \\
\midrule
All-trial accuracy & 0.525 & 0.435 & +0.090 \\
Matched-coverage accuracy & 0.613 & 0.524 & +0.089 \\
ECE & 0.141 & 0.176 & -0.035 \\
Brier score & 0.591 & 0.672 & -0.081 \\
\bottomrule
\end{tabular}
\end{center}

Lower ECE and Brier score indicate better probabilistic reliability. The
ensemble improves both metrics on the primary dataset while also improving
all-trial and matched-coverage accuracy. The Wilcoxon signed-rank test for
matched-coverage accuracy gives \(p = 0.037\) for the one-sided hypothesis
that the ensemble improves over LDA, with the important limitation that
\(n = 9\) subjects is small.

\subsection{MOABB External Validation Findings}

\begin{center}
\begin{tabular}{lrrrr}
\toprule
Dataset & Subjects & \(\Delta\) Acc. & \(\Delta\) ECE & \(\Delta\) Brier \\
\midrule
Cho2017 & 1 & +0.048 & +0.011 & -0.030 \\
PhysionetMI & 1 & +0.222 & +0.068 & -0.103 \\
\bottomrule
\end{tabular}
\end{center}

The MOABB smoke subset shows that the matched-coverage accuracy and Brier
score gains persist beyond BCI IV 2a. However, ECE is worse for the ensemble
on both one-subject external checks. This is a useful limitation: averaging
calibrated base probabilities can improve discrimination and squared
probability error while still requiring ensemble-level recalibration for
confidence-bin alignment. The larger MOABB run should therefore be used to
decide whether an extra calibration layer is needed after soft voting.

\subsection{Role in the Final Paper}

This v1 writeup should be used as the classical reliability baseline, not as a
separate final submission. Its paper-facing role is to motivate three choices
that carry into v2: leakage-aware cross-session evaluation, calibrated
probabilities, and selective prediction. The broader validation ingredients
now live in v2, where local BCI Competition IIIa, five-subject MOABB runs,
BNCI2014\_001, uncertainty-score ablations, temperature-scaling ablations, and
controller replay are evaluated under one output schema.

The v1 limitation is clear and should remain visible: equal soft voting
improves discrimination and Brier score, but it does not solve calibration
reliability by itself. That limitation is exactly why the final paper should
emphasize post-ensemble calibration and controller-level replay in v2.

\subsection{Reproducible Artifacts}

The held-out BCI IV 2a evaluation writes:
\begin{itemize}
  \item \texttt{outputs/heldout\_session/subject\_results.csv},
  \item \texttt{outputs/heldout\_session/risk\_coverage.csv},
  \item \texttt{outputs/heldout\_session/reliability\_metrics.csv},
  \item \texttt{outputs/heldout\_session/reliability\_diagram\_bins.csv},
  \item \texttt{outputs/heldout\_session/reliability\_diagrams/}.
\end{itemize}
The MOABB validation writes the same artifact types under
\texttt{outputs/moabb\_external/}. EEGNet can be added to these outputs with
the \texttt{--include-eegnet} flag after installing PyTorch.

\section{Conclusion}

M.I.N.D. v1 reframes motor-imagery EEG decoding as a reliability problem
rather than a pure accuracy problem. On BCI Competition IV 2a, the calibrated
LDA+SVM ensemble improves accuracy, ECE, and Brier score over calibrated LDA.
MOABB smoke validation suggests that the accuracy and Brier-score gains can
persist outside the primary benchmark, while also revealing a calibration
limitation: the ensemble may need post-vote recalibration to improve ECE
consistently.

The strongest v1 claim is therefore measured and baseline-oriented: a
lightweight, interpretable, calibrated ensemble with selective prediction can
make BCI commands more reliable than forced-choice LDA, and the final v2 paper
extends that claim with stronger deep models, broader external validation, and
device-facing controller replay.

\end{document}
